
\chapter{Einleitung}

Schach ist ein traditionsreiches Brettspiel, welches logisches Denken und Konzentration fördert. Durch den zunehmenden Einsatz digitaler Medien werden klassische Spiele jedoch immer häufiger durch rein elektronische Lösungen ersetzt. Aus diesem Grund entstand die Idee, ein herkömmliches Schachbrett mit moderner Technik zu erweitern und dadurch ein interaktives Spielerlebnis zu schaffen.

Diese Diplomarbeit befasst sich mit der Entwicklung eines interaktiven, smarten Schachbretts, das klassische Spielmechanismen mit moderner Technik verbindet. Durch die Kombination aus Sensorik, lokaler und externer Datenverarbeitung mittels passender Software soll ein System entstehen, das Spielzüge automatisch erkennt, den Spieler aktiv unterstützt und nach der Partie zeigt welche Züge besser geeignet gewesen wären. Dadurch wird das Schachspiel sowohl für Anfänger als auch für erfahrene Spieler attraktiver gestaltet.

Die Aufgabenstellung umfasst die Planung und Umsetzung der elektronischen und softwaretechnischen Komponenten sowie deren Integration in ein funktionales Gesamtsystem. Das Projekt verbindet technische Fachbereiche wie die Entwicklung und Programmierung der benötigten Hardware und Informationsaustausch mit einem speziellen Server auf dem der Nutzer seine Partien analysieren kann.

\section{Vertiefende Aufgabenstellung}

\subsection{Schwaninger Christoph}

\subsubsection{Bauteilauswahl}

Auswahl aller Hardware-Komponenten, welche für die Umsetzung des intelligenten Schachbretts benötigt werden.

\subsubsection{Schaltungsentwicklung}

Entwicklung einer Verschaltungslogik welche die Ansteuerung und Informationsausgabe der Felder effizient ermöglicht.

\subsubsection{Platinendesign}

Die Zuvor entwickelte Verschaltung der Komponenten für einen Platinenfertiger nach Anforderung erstellen.

\subsubsection{Gehäusefertigung}

Design, assemblieren und Einbindung aller benötigten Komponenten in das Gehäuse.

\subsubsection{Stromversorgung}

Auswahl und Verkabelung eines Stromversorgungssystem, welche den Anforderung aller Hardware-Komponenten entspricht.


\subsection{Koch Maximilian}

\subsubsection{Steuerung einer 2x2 Matrix}
Bei näherung eines Magneten an einen Sensor, wird dessen dasein vom Controller erkannt.

\subsubsection{Rotary Encoder als Eingabegerät}
Anhand einer vorher selbstgebauten Platine die Position von Magneten in einer Matrix feststellen.

\subsubsection{Entwicklung eines Menüs}
Entwicklung eines Menüs auf dem OLED-Bildschirm des Schachbretts. Eine Auswahl zwischen Spieler gegen Spieler oder Spieler gegen Computer erstellen.

\subsubsection{Einbindung der Spielererkennung}
Durch das Auslesen des eingebauten RFID Sensors die Unterscheidung verschiedener Spieler sicherstellen.


\subsubsection{Die Spiellogik}
Erstellen einer Spiellogik, mit der man sowohl gegen einen Schachcomputer als auch gegen einen menschlichen Gegner 
spielen kann.

\subsubsection{Die Datenübertragung}
Über Sockets eine Verbindung zu einem Server herstellen, um die gespielte Partie zu übertragen.


\subsection{Seidner Christina}

\subsubsection{Die Website}
Auf der Website können sich Spieler mithilfe des NFC-Sensors am Schachbrett registrieren, um dann ihre gespielten Spiele zu analysieren. Zudem kann man auf der Website Schach spielen. 

\subsubsection{Die Datenbank}
Auf der Datenbank werden alle relevanten Daten gespeichert, wie beispielsweise Benutzername, Rückgabewert des NFC-Sensors und die Schachzüge. 

\subsubsection{Der Server}
Der Server verbindet Website, Datenbank und Schachbrett miteinander. 

\newpage

\section{Dokumentation der Arbeit}

Es werden die Projektergebnisse dokumentiert

\begin{itemize}
	\item Grundkonzept
	\item Theoretische Grundlagen
	\item Praktische Umsetzung
	\item Lösungsweg
	\item Alternativer Lösungsweg
	\item Ergebnisse inkl. Interpretation
\end{itemize}

Weitere Anregungen:

\begin{itemize}
	\item Fertigungsunterlagen
	\item Testfälle (Messergebnisse…)
	\item Benutzerdokumentation
	\item Verwendete Technologien und Entwicklungswerkzeuge
\end{itemize}